%% ================================================================
%% Artigo Cienciométrico -- Ligas de Alta Entropia para Aplicações Hipersônicas
%% Configurações e Preâmbulo -- Padrão Moderno Overleaf / Academic Template
%% ================================================================
\documentclass[11pt,a4paper]{article}

%% ---------------- Codificação e Fontes
\usepackage[utf8]{inputenc}
\usepackage[T1]{fontenc}
\usepackage{textcomp}
\usepackage{lmodern}

%% ---------------- Paleta de Cores Acadêmica Moderna (Overleaf Style)
\usepackage{color}
\definecolor{primary}{rgb}{0.07, 0.21, 0.42}       % Deep Navy / Azul Marinho Nobre
\definecolor{secondary}{rgb}{0.18, 0.40, 0.62}     % Slate Blue / Azul Médio
\definecolor{accent}{rgb}{0.65, 0.15, 0.15}        % Dark Crimson / Destaque sutil
\definecolor{linegray}{rgb}{0.80, 0.84, 0.90}      % Linha divisória suave
\definecolor{boxbg}{rgb}{0.96, 0.97, 0.99}        % Fundo suave para caixas de resumo
\definecolor{tablehead}{rgb}{0.90, 0.93, 0.97}     % Fundo para cabeçalho de tabelas

%% ---------------- Geometria e Espaçamento de Página
\usepackage[top=2.5cm,bottom=2.5cm,left=2.5cm,right=2.5cm,headheight=14.5pt]{geometry}
\renewcommand{\baselinestretch}{1.18}               % Entrelinha arejada e agradável à leitura
\setlength{\parindent}{1.0cm}                       % Recuo padrão elegante
\setlength{\parskip}{3pt plus 1pt minus 1pt}        % Leve espaçamento entre parágrafos
\usepackage{indentfirst}

%% ---------------- Cabeçalhos e Rodapés Dinâmicos (fancyhdr)
\usepackage{fancyhdr}
\pagestyle{fancy}
\fancyhf{}
\fancyhead[L]{\small\textsf{\color{secondary} Estudo Cienciométrico --- Ligas de Alta Entropia para Voo Hipersônico}}
\fancyhead[R]{\small\textsf{\bfseries\color{primary} \thepage}}
\renewcommand{\headrulewidth}{0.5pt}
\renewcommand{\headrule}{{\color{linegray}\hrule height \headrulewidth \vspace{2pt}}}
\renewcommand{\footrulewidth}{0pt}

%% Estilo limpo para primeira página
\fancypagestyle{plain}{%
  \fancyhf{}%
  \fancyfoot[C]{\small\textsf{\color{secondary} \thepage}}%
  \renewcommand{\headrulewidth}{0pt}%
}

%% ---------------- Nomes em Português
\renewcommand{\abstractname}{Resumo}
\renewcommand{\figurename}{Figura}
\renewcommand{\tablename}{Tabela}
\renewcommand{\refname}{Referências Bibliográficas}

%% ---------------- Matemática, Gráficos e Tabelas
\usepackage{amsmath,amssymb}
\usepackage{graphicx}
\usepackage{tabularx}
\usepackage{array}
\usepackage{multicol}

%% ---------------- Macros de compatibilidade e elegância para tabelas
\providecommand{\toprule}{\hline\hline}
\providecommand{\midrule}{\hline}
\providecommand{\bottomrule}{\hline\hline}

%% ---------------- Legenda de Figuras e Tabelas
\newcommand{\legend}[1]{\par\vspace{4pt}{\footnotesize\color{secondary}\centering\textsf{#1}\par}\vspace{6pt}}

%% ---------------- Citação longa em bloco
\newenvironment{citacao}
  {\par\vspace{0.6em}\begin{list}{}{\setlength{\leftmargin}{3.5cm}\setlength{\rightmargin}{0.5cm}\setlength{\parsep}{0pt}\setlength{\topsep}{0pt}}
   \item[]\small\color{primary}\renewcommand{\baselinestretch}{1.05}\selectfont}
  {\end{list}\vspace{0.6em}}

%% ---------------- Tipografia de Títulos (Estilo Editorial Overleaf)
\makeatletter
\renewcommand{\section}{\@startsection{section}{1}{\z@}%
  {-3.5ex \@plus -1ex \@minus -.2ex}%
  {1.8ex \@plus.2ex}%
  {\normalfont\Large\bfseries\color{primary}}}

\renewcommand{\subsection}{\@startsection{subsection}{2}{\z@}%
  {-2.8ex\@plus -1ex \@minus -.2ex}%
  {1.2ex \@plus .2ex}%
  {\normalfont\large\bfseries\color{secondary}}}

\renewcommand{\subsubsection}{\@startsection{subsubsection}{3}{\z@}%
  {-2.2ex\@plus -1ex \@minus -.2ex}%
  {0.8ex \@plus .2ex}%
  {\normalfont\normalsize\bfseries\color{primary}}}
\makeatother

%% ---------------- Links e URLs formatadas
\setlength{\leftmargini}{1.0cm}
\usepackage{url}
\urlstyle{sf}
\providecommand{\href}[2]{{\color{secondary}\underline{#2}}}

%% ---------------- Título e Autoria Estilizados
\title{\vspace{-1.5cm}%
  {\Huge\bfseries\color{primary} Estudo Cienciométrico Avançado sobre Ligas de Alta Entropia para Aplicações Hipersónicas}\\[0.35cm]
  {\Large\textsl{\color{secondary} Redes de Autoria, Artigos Seminais e Evolução do Conhecimento}}%
}
\author{%
  \textbf{Valber Marcelino Filho}\thanks{Instituto Militar de Engenharia / Pesquisador Principal} \quad\and\quad
  \textbf{Rodrigo Andrade Paes}\thanks{Orientador e Supervisor de Pesquisa}%
}
\date{\vspace{0.2cm}\small\textsf{\color{secondary} Outubro de 2026}}

%% ---------------- Formatação refinada de tabelas
\renewcommand{\arraystretch}{1.28}
\setlength{\tabcolsep}{6pt}
